%% Chapter 4 — System Integration
\section{System Integration}
\label{sec:integration}

\subsection{FPGA Configuration and Board Bring-up}

The ZCU216 boots the PYNQ Linux image from an SD card.  After boot, the
\QICK{} bitstream is loaded onto the FPGA fabric by instantiating a
\texttt{QickSoc} object in Python, which calls the PYNQ overlay loader and
configures all IP blocks.  Because the RFSoC converters are controlled through
AXI (Advanced eXtensible Interface) register maps rather than external
chip-select lines, initialisation is
entirely software-driven.

The board is accessed over SSH, and experiments are run from Jupyter notebooks
served from the on-board PYNQ environment.  A single \texttt{QickSoc} instance,
created once per session, loads the bitstream and exposes the converters and
tProcessor to the notebook; the phase calibration of
Section~\ref{subsec:phase_calib} then holds for the remainder of that session.

\subsection{Channel Mapping}

The \QICK{} firmware assigns each physical DAC output to a numbered
\emph{generator channel} and each physical ADC input to a numbered
\emph{readout channel}.  This assignment is specific to the firmware build:
different bitstreams expose different channel counts and even renumber the same
physical port.  The experiments here used the tProcessor v2 firmware build
\texttt{2024-09-03\_216\_tprocv2r21\_demo}~\cite{QICKtprocv2}; Table~\ref{tab:channels}
lists the generator and readout channels used, with their physical ports.

\begin{table}[ht]
\centering
\caption{QICK channel mapping used on the ZCU216 (tProcessor v2 build).}
\label{tab:channels}
\begin{tabular}{llll}
\toprule
\QICK{} Channel & Physical Port & HD Header & Function \\
\midrule
Generator 0 (\texttt{axis\_signal\_gen\_v6}) & DAC tile 2, blk 0 & JHC3 & Readout drive \\
Readout 0                                    & ADC tile 2, blk 0 & JHC7 & IQ acquisition \\
\bottomrule
\end{tabular}
\end{table}

A second path (generator 3, the \texttt{axis\_sg\_mux8\_v1} multiplexed engine, into
readout 6 on JHC8) was also exercised, to compare a filtered against an unfiltered
front end; the single \texttt{axis\_signal\_gen\_v6} path above is the one used for
the measurements reported here.

\subsection{Phase Calibration}
\label{subsec:phase_calib}

Each time the ZCU216 is powered on or the FPGA bitstream is reloaded, the DDS
oscillators in the DAC and ADC blocks boot with a random mutual phase offset
$\phi$.  Without compensation, IQ measurements rotate unpredictably between
sessions, making it impossible to compare results across power cycles.

\paragraph{Phase model.}
A loopback tone from the readout DAC to the readout ADC measures $\phi$ at a given
frequency.  Because the two oscillators are effectively separated by a fixed time
delay $\tau$ (the total group delay around the loopback path), the phase
offset scales linearly with frequency:
\begin{equation}
  \phi(f) = \phi_0 - 360^\circ \cdot \tau \cdot f,
\label{eq:phase_model}
\end{equation}
where $\phi_0$ is a constant offset and $\tau$ is extracted from the fit.
The fitted $\tau$ reflects the combined pipeline delay of the DAC and ADC signal
chains, including the digital up- and down-conversion stages (DUC/DDC) and internal
FIFOs, rather than the physical cable propagation delay, which for a bench loopback
of 1--2\,m contributes only a few nanoseconds.

\paragraph{Calibration procedure.}
The phase calibration sweeps a range of frequencies $f_i$ near the target readout
frequency, records the averaged IQ phasor at each point, and extracts
$\phi_i = \arg\langle I + jQ \rangle$.  The pair $(\phi_0, \tau)$ is then fitted
using \texttt{scipy.optimize.least\_squares}, taking care to handle the
$0^\circ/360^\circ$ phase wrap correctly.  The fit is considered reliable when
the signal-to-noise ratio satisfies $|\text{mag}| \gg \text{RMS}$ (e.g.
$637 \gg 1.8$~ADU in a typical calibration run), where ADU denotes analog-to-digital
units, the raw integer output of the ADC.

The resulting \texttt{res\_phase} value:
\begin{equation}
  \phi_\text{res} = \phi_0 - 360^\circ \cdot \tau \cdot f_\text{res}
\end{equation}
is stored in the experiment configuration dictionary and passed to the readout
pulse via \texttt{soccfg.deg2reg(res\_phase)}, which converts degrees to the
integer register format expected by the tProcessor.  This calibration must be
re-run after every board reboot or bitstream reload, but remains valid throughout
a measurement session.

\paragraph{Measurement cadence.}
A summary of the recommended start-of-session checklist is:
\begin{enumerate}
  \item Power on and wait for PYNQ to boot ($\approx$\SI{60}{\second}).
  \item Connect the loopback (readout DAC $\to$ balun $\to$ SMA $\to$ balun $\to$
        readout ADC) and run the phase calibration notebook.
  \item Identify the correct \texttt{adc\_trig\_offset} using decimated readout,
        then switch to accumulated readout for all data-taking
        (Section~\ref{subsec:decimated_vs_accum}).
  \item At session end, call \texttt{soc.reset\_gens()} to halt all running
        generators.
\end{enumerate}
